\section{Canonical harmonic coordinates}\label{sec:harmonic}

We use the full endomorphism complex until imposing the determinant
equation. At the HE background, the holomorphic zero-forms are parallel,
so
\[
 P=H^0(X,\End E_0)=\bigoplus_i I_{F_i}\otimes\End W_i.
\]
Its trace-free part is identified with $\g_0$ in
\eqref{eq:finite-spectrum}. Write $Q_0$ for the orthogonal complement
of $P_0=P\cap\End_0E_0$ in trace-free zero-forms, and $Q_H$ for the
corresponding projection on Hermitian trace-free zero-forms.
All complements in this section are fixed at $(h_0,\omega_0)$.
Choose an integer $k>n+2$ and use $W^{k,2}$ for one-forms and
$W^{k+1,2}$ for complex gauges.

\subsection{Convergent reconstruction and exact equivariance}

\begin{lemma}\label{lem:kuranishi}
The recursion \eqref{eq:recursion-intro} has a unique analytic solution
on a small ball in $V$, smooth in the spatial variables. Writing
$r=\|x\|$ and $J_x=D_xa$, one has, in every fixed $C^j$ norm,
\begin{equation}\label{eq:canonical-estimates}
 \begin{gathered}
 a=x+O(r^2),\qquad J_x=\iota+O(r),\\
 H_1a=x,\qquad d^*a=0,\qquad H_1J_x=I,\qquad d^*J_x=0 .
 \end{gathered}
\end{equation}
Here $\iota$ is harmonic inclusion. The reconstructed operator is
integrable exactly when $\kappa(x)=0$, and $\tr a=\tr x$.
The maps are locally $G$-equivariant. With
$D_a=d+[a,\cdot]$ on endomorphism forms,
\begin{equation}\label{eq:exact-action}
 D_a\xi=-J_xC_x\xi\qquad(\xi\in P).
\end{equation}
The identity holds on the whole ambient ball.
\end{lemma}

\begin{proof}
Sobolev multiplication and the bounded map
$d^*G_2:W^{k,2}\to W^{k+1,2}$ make the right side of
\eqref{eq:recursion-intro} a contraction on a small ball.
Analytic parameter differentiation gives the first two estimates.
The nonlinear term is coexact, proving the harmonic and co-closed
identities. Bootstrapping the equation gives spatial smoothness and
the estimates in any specified finite regularity order.

Put $\mathcal R=da+a\wedge a$. Hodge decomposition gives
\[
 \mathcal R=H_2(a\wedge a)+d^*dG_2(a\wedge a),\qquad
 d^*\mathcal R=0,\qquad H_2\mathcal R=\kappa.
\]
The Bianchi identity is $d\mathcal R+[a,\mathcal R]=0$.
When $\kappa=0$, Hodge decomposition of $\mathcal R$ therefore implies
\[
 \mathcal R=d^*G_3d\mathcal R,\qquad
 \|\mathcal R\|_{W^{k-1,2}}\le Cr\|\mathcal R\|_{W^{k-1,2}}.
\]
For sufficiently small $r$ this forces $\mathcal R=0$.
In complex dimension less than three the degree-three term is absent,
with the same conclusion. Conversely, integrability implies $\kappa=0$
by harmonic projection. Since $\tr(a\wedge a)=0$ and trace commutes
with the scalar Hodge operators, $\tr a=\tr x$.

Every element of $G$ is parallel. Its action commutes with the Hodge
operators, even though the action need not be unitary. Uniqueness of
the contraction gives $a(gx)=ga(x)g^{-1}$ whenever the local expressions
are defined. Differentiation gives
$J_x[\xi,x]=[\xi,a]=-D_a\xi$, proving \eqref{eq:exact-action}.
\end{proof}

In particular the entire scalar harmonic deformation is removed by
$V_{\det}$, also when $H^1(X,\OO_X)\ne0$.
This subspace need not itself be a full quiver representation.
We apply finite balancing on the full representation $V$ and impose
$V_{\det}$ as an additional invariant analytic equation.
Neither the construction nor any later estimate replaces $\kappa$
by its quadratic term.

\subsection{The complementary Hermitian--Einstein equation}

On the fixed Hermitian bundle define the curvature moment with fixed
density $dV_0=\omega_0^n/n!$ by
\begin{equation}\label{eq:curvature-moment}
 \M_t(A)=
 \frac{iF_A^0\wedge\omega_t^{n-1}}{(n-1)!\,dV_0}.
\end{equation}
The superscript denotes trace-free curvature. A possibly nonintegrable
$(0,1)$-operator determines a unitary connection by its adjoint
$(1,0)$-part, so the following equation is defined on the ambient ball.

\begin{lemma}\label{lem:complement}
There is a unique small
$\sigma(t,x)\in Q_HW^{k+1,2}(\Herm_0E_0)$ satisfying
\begin{equation}\label{eq:complement-equation}
 Q_H\M_t\bigl(e^{\sigma(t,x)}\cdot(\ddbar_0+a(x))\bigr)=0.
\end{equation}
It is smooth and $K$-equivariant, and
\begin{equation}\label{eq:sigma-estimate}
 \|\sigma(t,x)\|_{C^j}\le C_j(|t|+r^2),\qquad
 \|D_x\sigma(0,x)\|_{V\to C^j}\le C_jr .
\end{equation}
\end{lemma}

\begin{proof}
At $(t,x,s)=(0,0,0)$ the derivative in $s$ of the left side is a
fixed nonzero multiple of the endomorphism Dolbeault Laplacian on
$Q_H$, from $W^{k+1,2}$ to $W^{k-1,2}$.
It is invertible by the fixed complementary elliptic estimate.
The parameter implicit-function theorem gives $\sigma$.
The linear $x$ derivative vanishes at the origin: the linear curvature
variation is contracted to an expression in $d^*x$ and its adjoint
by the K\"ahler identities, and $x$ is harmonic.
Taylor's theorem gives \eqref{eq:sigma-estimate}; elliptic regularity
gives the indicated spatial norms. Uniqueness gives $K$-equivariance.
\end{proof}

Set
\[
 \Psi_t(x)=e^{\sigma(t,x)}\cdot(\ddbar_0+a(x)),\qquad
 H_{t,x}=h_0e^{2\sigma(t,x)} .
\]
The latter is the transported metric on the original holomorphic
presentation when $x\in Z$. The determinant equation is separate:
multiplication of a bundle metric by $b_t$ solves it and leaves the
induced endomorphism metric unchanged.

For a Hermitian multiplicity tuple $S=(S_i)$ write
\begin{equation}\label{eq:tracefree-lift}
 S^0=\bigoplus_i I_{F_i}\otimes S_i-c(S)I,\qquad
 c(S)=\frac{\sum_i r_i\tr S_i}{R}.
\end{equation}
Define the finite covector by
\begin{equation}\label{eq:finite-curvature}
 \langle m_t(x),S\rangle
   =\int_X\tr\bigl(S^0\M_t(\Psi_t(x))\bigr)dV_0 .
\end{equation}
At $x=0$, equivariance makes $\sigma(t,0)$ scalar on each multiplicity
space, preserving the factor decomposition. Chern--Weil on each
factor and the lift \eqref{eq:tracefree-lift} give
\begin{equation}\label{eq:central-value}
 m_t(0)=-\zeta(t),\qquad \rho_t=m_t-m_t(0).
\end{equation}
The pairing in \eqref{eq:finite-curvature} is the ordinary
multiplicity trace, not the spectral mass in \eqref{eq:finite-spectrum}.

\subsection{The quadratic ambient estimate}

Use the Hermitian $L^2$ pairing linear in the first variable, and
$\omega_{\mathrm{flat}}(v,w)=\operatorname{Im}\langle v,w\rangle$.
The affine space of unitary connections has its constant symplectic
form determined by $\omega_t$. Its pullback under $\Psi_t$ is denoted
$\Omega_t$, with the normalization
\begin{equation}\label{eq:moment-identity}
 d\langle\rho_t,S\rangle=2\Omega_t(iC_xS,\cdot),\qquad \rho_t(0)=0 .
\end{equation}
This is the gauge moment identity, using the $K$-equivariance of
$\Psi_t$. The form is closed and invariant on the ambient ball;
it need not have type $(1,1)$ for the finite complex structure.

\begin{proposition}\label{prop:quadratic-form}
In the canonical harmonic coordinates,
\begin{equation}\label{eq:quadratic-form}
 \|\Omega_{t,x}-\omega_{\mathrm{flat}}\|
                 \le C(|t|+\|x\|^2)
\end{equation}
as real bilinear forms on the entire ambient representation.
\end{proposition}

\begin{proof}
It suffices to prove vanishing of the spatial first derivative at
$(t,x)=(0,0)$. Write the homogeneous quadratic terms as
\[
 a(x)=x+a_2(x)+O(r^3),\qquad
 \sigma(0,x)=\sigma_2(x)+O(r^3).
\]
Here
\[
 a_2(x)=-d^*G_2(x\wedge x)\in\im d^* .
\]
For the action $g\cdot\ddbar=g\ddbar g^{-1}$, the $(0,1)$ component is
\begin{equation}\label{eq:two-jets}
 \Psi_0(x)-\ddbar_0
        =x+a_2(x)-d\sigma_2(x)+O_{C^j}(r^3).
\end{equation}
Both displayed quadratic terms are orthogonal to harmonic one-forms.
Their real parameter derivatives are also orthogonal to the entire
harmonic space. Every term linear in $x$ in the pullback of the
constant affine form pairs a harmonic tangent with such a derivative,
so vanishes. The remaining error is $O(r^2)$ in operator norm.
Smooth dependence on $t$ contributes $O(|t|)$.
\end{proof}

This is a statement about the canonical recursion, not all analytic
coordinates tangent to the identity. A change $x\mapsto x+F_2(x)+\cdots$
can introduce a harmonic quadratic term, leaving only an $O(r)$
estimate in the new coordinates. The distinction will matter for the
formality coordinates in Section~\ref{sec:recovery}.

\subsection{A positive form before the finite quotient}

For $x,v\in V_{\det}$ define
\begin{equation}\label{eq:complement-minimum}
 \B_{t,x}(v,v)=\inf_{\eta\in Q_0W^{1,2}}
       \|J_xv-D_a\eta\|^2_{H_{t,x},\omega_t}.
\end{equation}
The fixed estimate $\|\eta\|_{W^{1,2}}\le C\|d\eta\|$ on $Q_0$
and multiplication by $a=O(r)$ give, by absorption,
\begin{equation}\label{eq:complement-coercive}
 \|\eta\|_{W^{1,2}}\le C\|D_a\eta\|_{H_{t,x},\omega_t}.
\end{equation}
Thus the normal equation for this minimum has a uniformly coercive
operator on the fixed Hilbert space $Q_0W^{1,2}$.

\begin{proposition}\label{prop:positive-form}
The form $\B_{t,x}$ is smooth, Hermitian and positive on a common
neighborhood, and
\begin{equation}\label{eq:B-comparison}
 (1-C\delta)\|v\|^2\le\B_{t,x}(v,v)\le(1+C\delta)\|v\|^2,
 \qquad \delta=|t|+r^2.
\end{equation}
For a genuine integrable tangent $v$ the full gauge minimum is
\begin{equation}\label{eq:full-minimum}
 \inf_{\gamma\in W^{1,2}(\End_0E_0)}
       \|J_xv-D_a\gamma\|^2
 =\inf_{\xi\in\g_0}\B_{t,x}(v-C_x\xi,v-C_x\xi).
\end{equation}
All norms in the minima use $H_{t,x},\omega_t$.
\end{proposition}

\begin{proof}
First let $t=0$. Since $H_1J_x=I$, harmonic orthogonality gives
\[
 \|J_xv\|^2=\|v\|^2+O(r^2)\|v\|^2 .
\]
The metric $H_{0,x}$ changes squared endomorphism norms by
$1+O(r^2)$. In the reference norm, $d^*J_xv=0$, so
\[
 |\langle J_xv,D_a\eta\rangle_{H_{0,x},\omega_0}|
             \le Cr\|v\|\|\eta\|_{W^{1,2}}.
\]
Indeed the $d\eta$ term is zero before metric variation, the
$[a,\eta]$ term is $O(r)$, and the metric variation contributes
$O(r^2)$. Coercivity bounds the minimizer by $Cr\|v\|$ in
$W^{1,2}$. The removed energy is at most $Cr^2\|v\|^2$.
This proves \eqref{eq:B-comparison} at $t=0$. Smooth dependence
of the uniformly coercive inverse contributes $O(|t|)$.
The inverse also proves smoothness of $\B$.

Every gauge zero-form splits uniquely as $\xi+\eta$, with
$\xi\in P_0$ and $\eta\in Q_0W^{1,2}$. Substitute
\eqref{eq:exact-action}, minimize over $\eta$, and replace $\xi$ by
$-\xi$ to obtain \eqref{eq:full-minimum}.
For an integrable tangent, differentiating the full integrability
equation gives $D_aJ_xv=0$. Hodge minimization therefore identifies
the left side with the harmonic Kodaira--Spencer squared norm.
\end{proof}

The same construction on all of $V$ uses the full complement of $P$,
including gapped scalar zero-forms. On $V_{\det}$ it agrees with
\eqref{eq:complement-minimum}, since trace and trace-free forms are
orthogonal and preserved by $D_a$. No inverse of $C_x^*C_x$ or
$D\kappa_x$ enters the proof.

